%Two resources useful for abstract writing.
% Guidance of how to write an abstract/summary provided by Nature: https://cbs.umn.edu/sites/cbs.umn.edu/files/public/downloads/Annotated_Nature_abstract.pdf %https://writingcenter.gmu.edu/guides/writing-an-abstract
\chapter*{\center \Large  Abstract}
%%%%%%%%%%%%%%%%%%%%%%%%%%%%%%%%%%%%%%
% Replace all text with your text
%%%%%%%%%%%%%%%%%%%%%%%%%%%%%%%%%%%
Packet-processing applications in middleboxes such as firewalls and load
balancers are commonly written as sequential C programs that are typically loopless and parse incoming
packets by extracting header fields and branching on their values. However, it is hard to build automatic optimisers directly from this representation.


This project presents a synthesis-based compiler that automatically
translates sequential C-like packet parsing code into an equivalent
finite-state machine (FSM). The compiler involves a
Counterexample-Guided Inductive Synthesis (CEGIS) loop using the Z3 SMT
solver, and produces a tabular FSM that maps to P4,
a language for programmable network hardware. The CEGIS verifier checks all possible input packets to confirm that the synthesised FSM produces the same \textsc{Accept}/\textsc{Reject} verdict and extracts the same set of header fields as the original code.


The FSM representation may be used for the following optimisations.

First, \textbf{payload parking}: The generated P4 code can be used to 
program hardware to extract and forward only the relevant headers,
avoiding unnecessary movement of packet payloads through middleboxes that
act only on header fields. 

Second, \textbf{batch processing}: The synthesised FSM captures the control flow structure that arises due to parsing. When no other branches exist in the processing code, packets in the same parse state follow identical execution until the next branch point in the parse graph and can therefore be grouped and processed together using vector instructions. Even when additional branches exist (for example, due to TTL checks or lookup-dependent logic), further analyses may be done to handle divergence due to factors other than parsing.

The FSM thus provides a foundation for further analyses targeting hardware-accelerated packet processing.



% \vspace{2em}
% This is a project report template, including instructions on how to write a report. It also has some useful examples to use \LaTeX. Do read this template carefully. The number of chapters and their titles may vary depending on the type of project and personal preference. Section titles in this template are illustrative should be updated accordingly. For example, sections named ``A section...'' and ``Example of ...'' should be updated. The number of sections in each chapter may also vary. This template may or may not suit your project. Discuss the structure of your report with your supervisor.

% %%%%%%%%%%%%%%%%%%%%%%%%%%%%%%%%%%%%%%%%%%%%%%%%%%%%%%%%%%%%%%%%%%%%%%%%%s
% ~\\[1cm]
% \noindent % Provide your key words
% \textbf{Keywords:} a maximum of five keywords/keyphrase separated by commas

% \vfill
% \noindent
% \textbf{Report's total word count: Following the abstract, the word count must be stated.} We expect at least 10,000 words in length and at most 15,000 words (starting from Chapter 1 and finishing at the end of the conclusions chapter, excluding references, appendices, abstract, text in figures, tables, listings, and captions), about 40 - 50 pages. \newline
% \newline
% \noindent
% \textbf{Program code should be uploaded to gitlab, and the gitlab link should be included alongside the word count, following the abstract.} \newline
% \newline
% You must submit your dissertation report (preferred in a PDF file) via the “Turnitin assignment” in Blackboard Learn by the deadline. If a student has resits from the taught modules, the dissertation deadline will be extended for 3 weeks from the original dissertation deadline.

